
% Java Style lstlisting

\lstdefinestyle{javaStyle}{
            language=Java,
            backgroundcolor=\color{gray!5},
            basicstyle=\ttfamily\small,
            keywordstyle=\color{blue!80!black}\bfseries,
            stringstyle=\color{teal!70!black},
            commentstyle=\color{gray!80}\itshape,
            numbers=left,
            numberstyle=\tiny\color{gray},
            stepnumber=1,
            frame=single,
            framerule=0.5pt,
            breaklines=true,
            showstringspaces=false,
            tabsize=2,
            morekeywords={var},
            xleftmargin=0pt,
            framexleftmargin=0pt
        }

\lstnewenvironment{javaSnippet}{
  \lstset{style=javaStyle}
}{}


% New list Style, i. ii. iii. ...
\newlist{enumRoman}{enumerate}{1}
\setlist[enumRoman]{label=\roman*.}

% New list Style, a) b) ...
\newlist{enumAlph}{enumerate}{1}
\setlist[enumAlph]{label=\alph*)}

% Java Console
    % Colors

\definecolor{consolebg}{RGB}{245,245,245}   % light gray background, like printed output
\definecolor{consoletext}{RGB}{0,0,0}      % black text

\renewcommand{\ttdefault}{lmtt} % Latin Modern Typewriter, looks nicer than default

    % Define style
    
\lstdefinestyle{javaOutputStyle}{
    basicstyle=\ttfamily\small\color{consoletext}, % monospaced text
    backgroundcolor=\color{consolebg},            % light gray background
    breaklines=true,
    showstringspaces=false,
    frame=none,                                   % NO horizontal lines
    xleftmargin=0pt,
    xrightmargin=0pt,
    aboveskip=1em,
    belowskip=1em,
}

\lstnewenvironment{javaOutput}{
    \lstset{style=javaOutputStyle}
}{}

\SetKwInOut{Input}{Input}
\SetKwInOut{Output}{Output}
